\section{Introduction}
\label{sec:intro}

The cost of a diffusion-transformer forward pass grows superlinearly with
resolution: token count grows quadratically with edge length, and
attention cost grows again in token count. That makes the title question
worth asking precisely. At high noise, can the \emph{gradient} of an
adapter be computed on a downscaled latent as a drop-in substitute for
the native one, with the same model, the same objective, and the same
$\sigma$ distribution, only the spatial grid changed on some steps?

There is a well-studied reason to hope the answer is yes. As the noise
level $\sigma$ grows, per-frequency signal-to-noise falls, and high
spatial frequencies drop below the noise floor first, so that above a
spectral crossover a downscaled latent carries \emph{almost} the full
surviving signal. But the studies of this resolution gap live almost
entirely on the \emph{inference} side: scale-wise distillation
\citep{starodubcev2025swd}, spectral progressive diffusion
\citep{xiao2026spd}, and spectrally-guided noise schedules
\citep{esteves2026spectral} run high-noise \emph{sampling} steps at
reduced resolution and judge the premise by the generated output. And
where resolution is reduced during training (pyramidal flow
matching, \citealp{jin2024pyramidal}; progressive-resolution training
more broadly,
\citealp{karras2018progressive,chen2024pixartalpha,hoogeboom2023simple}),
the low-resolution stage is given its own, modified objective, and
success is again read off the final samples. By the spectral premise,
substitution is safe \emph{whenever the noise has masked the
frequencies the coarse grid loses}.

None of this work, however, analyzes the gap at \emph{training} time,
that is, whether a downscaled step delivers the native gradient under
an unchanged objective; that is the route this paper takes. What a training step consumes is the expected adapter gradient,
and \emph{demotion} --- our term for the training-time substitution of
a downscaled, re-encoded latent for the native one --- perturbs more
than the network's input: the regression
target carries the clean image at unit weight at every noise level,
and the compute graph itself changes with the grid. Nor does
output-level evidence settle the matter: a small quality-score
difference does not transfer to a small difference in what a training
step \emph{learns}. This paper measures the question directly, at the
gradient level, holding model and objective fixed.

Measuring it honestly is, unexpectedly, a metrology problem. The
natural gap estimator (a redraw floor minus a single demoted-arm
cosine) is biased by finite-draw variance, and the bias grows as the
demoted grid shrinks: it \emph{manufactures} exactly the
token-count-ordered floor signature under test. Every claim in this
paper is therefore stated in debiased units, at the instrument's
measured resolution.

The contributions:

\begin{enumerate}
\item \textbf{Gap metrology and a two-term account of what demotion
costs} (\S\ref{sec:theory}). We define the
demotion gap at the gradient level and build a debiased instrument for
it: per-arm self-floors with a split-half attenuation correction, a
draw-count extrapolation $\gap(D) = \gap_\infty + c/D$, and a
certification resolution $\epsstar(N, D)$ under which ``safe'' is a
one-sided non-inferiority statement, not an absolute claim. From the
perturbation expansion and cosine geometry, the gap reduces under four
named assumptions to a \emph{two-term account}: a route amplitude on a
universal mismatch-amplitude curve read through the U-shaped gradient
norm, plus a $\sigma$-independent graph floor. The account connects directly to
the inference-side literature: the spectral argument, in its strongest
tolerance-parameterized form (SPD, \citealp{xiao2026spd}), yields its
own account of the data term's input-mediated part, one that is
ratio-only and floorless by construction.
\item \textbf{A measured safety map and a predictive tool}
(\S\ref{sec:evidence}, \S\ref{sec:practical}). Measured with that
instrument, the mildest route ($1024{\to}896$) is indistinguishable
from the re-encode control for $\sigma > 0.5$ at instrument
resolution, while no noise level rescues a $2\times$ downscale, and
no member of the spectral tolerance family reproduces the measured
map. The two-term account \emph{predicts} held-out routes at
${\sim}0.07$--$0.09$ RMSE, with the floor law correctly predicting
both held-out floors: a tool for anticipating a route's gap before
measuring it. Gating substitution on the measured-safe windows gives an opt-in
$\sigma$-conditional trainer realizing $-15.1\%$ forward FLOPs and
$-14.6\%$ wall-clock at a fixed step budget, with a sample-level
footprint within the training-seed lottery (the seed-noise yardstick,
Table~\ref{tab:yardstick}); generation-quality non-inferiority at
production scale remains the stated open gate.
\end{enumerate}

We first situate the question in the prior literature
(\S\ref{sec:related}); \S\ref{sec:theory} defines the estimand and its
certification resolution, derives the angular account of the demotion
gap and the gradient gap the spectral argument implies for its data
term, and constructs the debiased estimator; \S\ref{sec:evidence}
builds the instrument and discharges both accounts' predictions
against the measurements;
\S\ref{sec:practical} gives the resulting safety map and the
$\sigma$-conditional trainer.

\section{Related work}
\label{sec:related}

\paragraph{Flow matching and DiT.}
Flow matching \citep{lipman2023flow} trains a continuous-time generative
model by regressing a velocity field along prescribed probability paths
between noise and data; rectified flow \citep{liu2023rectified} takes the
paths to be straight lines, giving the linear noising
$z_\sigma = (1-\sigma)x + \sigma\epsilon$ and the constant regression
target $v = \epsilon - x$ used throughout this paper, and this
formulation is the training objective of current large text-to-image
systems \citep{esser2024scaling}. The Diffusion Transformer
\citep{peebles2023dit} replaces the U-Net denoiser with a plain
transformer over patchified latent tokens, so a spatial grid enters the
network only as a token sequence: token count grows quadratically
with edge length and attention cost quadratically again in token count,
the cost structure that makes resolution the expensive axis. Spatial
position is carried by rotary position embeddings \citep{su2024roformer},
whose per-band phases accumulate with grid coordinates; this
grid-dependent coordinate system is the one whose contribution to the
demotion floor \S\ref{sec:decomposition} isolates by intervention. Fine-tuning is
LoRA \citep{hu2022lora}: low-rank adapter updates on a frozen backbone,
the gradient regime all our probes measure.

\paragraph{Scale-wise and progressive-resolution diffusion.}
Progressive growing \citep{karras2018progressive} and resolution curricula
\citep{chen2024pixartalpha,hoogeboom2023simple} train at increasing
resolution as a schedule, each low-resolution stage carrying its own
stage-specific objective. A more recent family ties resolution to the
noise level itself, on a spectral premise. Under flow-matching noising,
each spatial-frequency band of the signal retains an effective
signal-to-noise ratio that falls as $\sigma$ grows, and the finest bands
sink below the noise floor first; the noise level above which a band is
masked is that band's \emph{crossover} $\sigma$. Above the crossover of
the bands a coarse grid cannot represent, a downscaled latent should
therefore carry almost the full surviving signal. This resolution gap
has so far been studied on the \emph{inference} side. Scale-wise
distillation (SwD) \citep{starodubcev2025swd} states the claim
directly (above the crossover, the noised downscaled latent is
close in distribution to the downscaled noised latent) and samples
early steps at reduced resolution. Spectral Progressive Diffusion (SPD)
\citep{xiao2026spd} gives the argument its sharpest available form: a
tolerance-parameterized per-frequency activation time under a diagonal
Gaussian spectral model, a criterion on the \emph{Bayes-optimal
predictor's} per-band output rather than merely on input SNR, from
which an inference-time resolution schedule follows.
Spectrally-guided noise schedules \citep{esteves2026spectral} apply the
same per-frequency reasoning to schedule design, deriving bounds on the
useful minimum and maximum noise levels from the image spectrum; here
the spectral account operates as the field's working model of what
noise masks. Where the premise does enter training, as in pyramidal flow
matching \citep{jin2024pyramidal}, the objective is restaged per
pyramid level rather than left native, and the premise is again
validated by the generated output. None of these analyzes what a
low-resolution step does to the \emph{training gradient} of an
unchanged objective; that is the question we measure, holding model
and objective fixed. We find that input-level closeness does not
transfer at any tolerance.

\paragraph{Positional extension.}
Position interpolation \citep{chen2023pi}, banded/frequency-selective
variants (YaRN, \citealp{peng2024yarn}), and diffusion-specific dynamic
position extrapolation \citep{issachar2025dype,zhao2026sigma} rescale
rotary coordinates to evaluate a model at unseen grid sizes. We use exact
positional interpolation as a \emph{causal intervention} (matching the
downscaled grid's relative phase geometry to native) to decompose the
resolution-sensitivity floor, and adopt the $\sigma$-gated band schedule of
\citet{zhao2026sigma} as a training-time refinement.

